% ==========================================
% Kapitel 4: Komplexe Wegintegrale
% ==========================================
\section{Komplexe Wegintegrale}

\begin{defi}{4.1 Stückweise Stetigkeit und Differenzierbarkeit}
Sei $[a,b] \subset \R$.
\begin{itemize}
    \item[i)] Eine Funktion $f: [a,b] \to \C$ heißt \textbf{stückweise stetig} ($f \in C_{st}^0([a,b]; \C)$), wenn es eine Zerlegung $a = t_0 < t_1 < \dots < t_n = b$ gibt, sodass $f$ auf jedem Teilintervall stetig ist und stetig auf die Ränder $[t_{j-1}, t_j]$ fortgesetzt werden kann.
    \item[ii)] $f$ heißt \textbf{stückweise stetig differenzierbar} ($f \in C_{st}^1([a,b]; \C)$), wenn $f$ auf ganz $[a,b]$ stetig ist und es eine Zerlegung gibt, sodass $f$ auf den Teilintervallen stetig differenzierbar ist. Die Ableitung ist beschränkt.
\end{itemize}
\end{defi}

\begin{defi}{4.3 Wege und Integrationswege}
Sei $[a,b] \subset \R$. Eine stetige Abbildung $\gamma: [a,b] \to \C$ heißt \textbf{Weg}. 
Ist $\gamma$ zusätzlich stückweise stetig differenzierbar ($\gamma \in C_{st}^1$), so heißt $\gamma$ ein \textbf{$C_{st}^1$-Weg} oder \textbf{Integrationsweg}. 
Ist $\gamma$ sogar stetig differenzierbar auf $[a,b]$, heißt es \textbf{$C^1$-Weg}.
\end{defi}

\begin{defi}{4.4 Eigenschaften von Wegen}
Für einen Weg $\gamma: [a,b] \to \C$ definiert man:
\begin{itemize}
    \item $\gamma(a) \in \C$ als \textbf{Anfangspunkt} und $\gamma(b) \in \C$ als \textbf{Endpunkt}.
    \item $Sp(\gamma) := \gamma([a,b]) = \{\gamma(t) : t \in [a,b]\}$ als \textbf{Spur} des Weges.
    \item $\gamma$ heißt \textbf{geschlossen}, wenn $\gamma(a) = \gamma(b)$.
\end{itemize}
$\gamma$ wird hierbei auch als Parametrisierung bezeichnet.
\end{defi}

\begin{tcolorbox}[colback=white, colframe=black!70, title={Bsp. 4.5 Standard-Wege}]
Wichtige Parametrisierungen für die Übung:
\begin{itemize}
    \item \textbf{Kreislinie} um $z_0$ mit Radius $r$: $\gamma(t) = z_0 + r e^{\I t}$ für $t \in [0, 2\pi]$.
    \item \textbf{Strecke} von $z_1$ nach $z_2$: $\gamma(t) = (1-t)z_1 + t z_2$ für $t \in [0, 1]$.
\end{itemize}
\end{tcolorbox}

\begin{defi}{4.6 Integral komplexwertiger Funktionen}
Sei $f \in C_{st}^0([a,b]; \C)$ mit der Zerlegung in Real- und Imaginärteil $f(t) = u(t) + \I v(t)$. Das Integral wird komponentenweise definiert:
\[
    \int_a^b f(t) \, dt := \int_a^b u(t) \, dt + \I \int_a^b v(t) \, dt
\]
\end{defi}

\begin{satz}{4.7 Eigenschaften des Integrals (Prop)}
Seien $f,g \in C_{st}^0([a,b]; \C)$ und $\alpha, \beta \in \C$. Dann gilt:
\begin{itemize}
    \item[i)] \textbf{Linearität:} $\int_a^b (\alpha f(t) + \beta g(t)) \, dt = \alpha \int_a^b f(t) \, dt + \beta \int_a^b g(t) \, dt$.
    \item[ii)] \textbf{Additivität:} $\int_a^b f(t) \, dt = \int_a^c f(t) \, dt + \int_c^b f(t) \, dt$ für alle $c \in [a,b]$.
    \item[iii)] \textbf{Betragsabschätzung:} $\left| \int_a^b f(t) \, dt \right| \le \int_a^b |f(t)| \, dt$.
\end{itemize}
\end{satz}

\begin{satz}{4.8 Hauptsatz der Differential- und Integralrechnung (HDI)}
Sei $f \in C_{st}^1([a,b]; \C)$. Dann gilt:
\[
    \int_a^b f'(t) \, dt = \Big[ f(t) \Big]_{t=a}^{t=b} = f(b) - f(a)
\]
\end{satz}

\begin{defi}{4.10 Komplexes Wegintegral}
Sei $\gamma: [a,b] \to \C$ ein Integrationsweg ($C_{st}^1$-Weg) und $f: Sp(\gamma) \to \C$ eine stetige Funktion. Dann heißt
\[
    \int_\gamma f(z) \, \dz := \int_a^b f(\gamma(t)) \cdot \gamma'(t) \, dt
\]
das \textbf{komplexe Wegintegral} (oder Kurvenintegral).
\end{defi}

\begin{rezept}{Satz 4.11 Das wichtigste Standard-Integral}
Sei $z_0 \in \C$ und $\gamma(t) = z_0 + r e^{\I t}$ mit $t \in [0, 2\pi]$ die Kreislinie. Dann gilt für Potenzen $n \in \Z$:
\[
    \oint_\gamma (z-z_0)^n \, \dz = \begin{cases} 
        0 & \text{für } n \in \Z \setminus \{-1\} \\ 
        2\pi \I & \text{für } n = -1 
    \end{cases}
\]
\textit{Skill-Hinweis aus der Vorlesung: Dieses Integral muss man aus dem Effeff ausrechnen bzw. anwenden können!}
\end{rezept}

\begin{defi}{4.13 Weglänge}
Sei $\gamma: [a,b] \to \C$ ein Weg. Die \textbf{Länge} des Weges ist definiert als das Supremum über alle polygonalen Approximationen:
\[
    L(\gamma) = \sup \left\{ \sum_{i=1}^m |\gamma(t_i) - \gamma(t_{i-1})| \;\middle|\; a = t_0 < t_1 < \dots < t_m = b \right\}
\]
\end{defi}

\begin{satz}{4.14 Berechnung der Weglänge (Lemma)}
Ist der Weg ein Integrationsweg, also $\gamma \in C_{st}^1([a,b]; \C)$, so vereinfacht sich die Längenberechnung zu:
\[
    L(\gamma) = \int_a^b |\gamma'(t)| \, dt < \infty
\]
\end{satz}

% --- Neu aus der Vorlesung: Parameter-Transformationen ---

\begin{satz}{4.15 Eigenschaften des Wegintegrals}
Sei $\gamma \in C_{st}^1([a,b]; \C)$, $f, g: Sp(\gamma) \to \C$ stetig und $\alpha, \beta \in \C$. Dann gilt:
\begin{itemize}
    \item[i)] \textbf{Linearität:} $\int_\gamma (\alpha f(z) + \beta g(z)) \, \dz = \alpha \int_\gamma f(z) \, \dz + \beta \int_\gamma g(z) \, \dz$.
    \item[ii)] \textbf{Standardabschätzung:} $\left| \int_\gamma f(z) \, \dz \right| \le \int_\gamma |f(z)| \, |dz| \le \max_{z \in Sp(\gamma)} |f(z)| \cdot L(\gamma)$.
\end{itemize}
\end{satz}

\begin{defi}{4.16 Zusammengesetzte Wege und entgegengesetzte Wege}
\begin{itemize}
    \item[i)] Seien $\gamma_1: [a,b] \to \C$ und $\gamma_2: [b,c] \to \C$ Wege mit $\gamma_1(b) = \gamma_2(b)$. Der \textbf{zusammengesetzte Weg} $\gamma = \gamma_1 \cdot \gamma_2$ ist definiert als:
    \[
    \gamma(t) = \begin{cases} \gamma_1(t) & t \in [a,b] \\ \gamma_2(t) & t \in [b,c] \end{cases}
    \]
    \item[ii)] Für einen Weg $\gamma: [a,b] \to \C$ ist der \textbf{entgegengesetzte Weg} $\overleftarrow{\gamma}: [a,b] \to \C$ definiert durch $\overleftarrow{\gamma}(t) = \gamma(a+b-t)$.
\end{itemize}
\end{defi}

\begin{satz}{4.17 Additivität des Integrals}
Sei $\gamma = \gamma_1 \cdot \gamma_2$ ein zusammengesetzter Weg. Dann gilt:
\[
    \int_\gamma f(z) \, \dz = \int_{\gamma_1} f(z) \, \dz + \int_{\gamma_2} f(z) \, \dz
\]
\end{satz}

\begin{tcolorbox}[colback=white, colframe=black!70, title={Bsp 4.18: Parametrisierung der Kreislinie $\partial B_r(z_0)$}]
Die Kreislinie kann unterschiedlich parametrisiert werden:
\begin{itemize}
    \item[i)] $\gamma_1(t) = z_0 + re^{2\pi\I t}, t \in [0,1]$
    \item[ii)] $\gamma_2(t) = z_0 + re^{\I t}, t \in [0,2\pi]$
    \item[iii)] $\gamma_3$ (stückweise, siehe Vorlesungsskript)
    \item[iv)] $\gamma_4(t) = z_0 + re^{2\pi\I t}, t \in[1/2, 3/2]$
    \item[v)] $\gamma_5(t) = z_0 + re^{-2\pi\I t}, t \in [0,1]$ (durchläuft Kreis im Uhrzeigersinn).
\end{itemize}
Es gilt: $\int_{\gamma_1} f(z) dz = \int_{\gamma_2} f(z) dz = \int_{\gamma_3} f(z) dz = \int_{\gamma_4} f(z) dz = -\int_{\gamma_5} f(z) dz$.
\end{tcolorbox}

\begin{defi}{4.19 Parameter-Transformation}
Seien $\gamma: [a,b] \to \C$ und $\tilde{\gamma}: [\tilde{a},\tilde{b}] \to \C$ Wege. Eine bijektive Abbildung $\varphi: [\tilde{a},\tilde{b}] \to [a,b]$ heißt \textbf{Parameter-Transformation} (oder $C_{st}^1$-Umparametrisierung), falls $\varphi \in C_{st}^1$ und $\varphi' > 0$ gilt. 
Gilt $\tilde{\gamma}(t) = \gamma(\varphi(t))$, so entsteht $\tilde{\gamma}$ durch Umparametrisierung aus $\gamma$.
\end{defi}

\begin{satz}{4.20 Invarianz und Richtung}
Sei $\tilde{\gamma}$ eine $C_{st}^1$-Umparametrisierung von $\gamma$. Dann gilt:
\begin{itemize}
    \item[i)] $Sp(\gamma) = Sp(\tilde{\gamma})$ und die End-/Anfangspunkte sind gleich.
    \item[ii)] $\int_{\tilde{\gamma}} f(z) \, \dz = \int_{\gamma} f(z) \, \dz$.
    \item[iii)] Für den entgegengesetzten Weg gilt: $\int_{\overleftarrow{\gamma}} f(z) \, \dz = -\int_{\gamma} f(z) \, \dz$.
\end{itemize}
\end{satz}

\begin{satz}{4.21 Vereinfachung der Notation}
Sei $\gamma$ ein geschlossener $C_{st}^1$-Weg, der injektiv auf $[a,b)$ ist und gegen den Uhrzeigersinn durchlaufen wird. Sei $\Gamma = Sp(\gamma)$. Dann schreibt man kurz $\int_{\Gamma} f(z) \, \dz := \int_{\gamma} f(z) \, \dz$. 
\end{satz}

\begin{satz}{4.22 Vertauschungssatz}
Sei $\gamma$ ein $C_{st}^1$-Weg und $f_n \to f$ gleichmäßig auf $Sp(\gamma)$. Dann gilt:
\[
    \lim_{n \to \infty} \int_\gamma f_n(z) \, \dz = \int_\gamma f(z) \, \dz
\]
\end{satz}